% =====================  SECTION VI  =====================
\section{Conclusion and Future Work}\label{sec:conclusion}

\subsection{Conclusion}\label{subsec:conclusion}
This paper set out to build an integrated, explainable, and independently verifiable
decision-support pipeline that carries a single user's raw financial inputs through risk
assessment, safety enforcement, explanation, portfolio allocation, and tamper-evident
logging---a combination the reviewed literature treats as five largely separate problems.
The pipeline delivers that integration end to end: it engineers three interpretable ratios
from five raw inputs, derives Low, Medium, and High risk segments by clustering (silhouette
0.370, stable at an adjusted Rand index of 0.868), trains a leakage-aware MLP that
reproduces the risk band at 91.0\% accuracy and 0.881 macro-F1, routes every prediction
through a conservative guardrail that can only escalate risk, explains each decision with a
model-adaptive feature-impact method, maps the reconciled band to a risk-aligned
BTC/ETH/USDC allocation, and records a SHA-256-hashed, verifiable receipt for every
recommendation.

Equally central to the conclusion is what the work does not claim. The evaluation
establishes that the conservative guardrail, not the learned model, determines most final
recommendations; that the proxy labels are arithmetically entangled with the classifier
inputs; that the MLP is not statistically better than logistic regression; and that the
portfolio simulation rests on fixed synthetic returns. The contribution is therefore an
engineering one---a coherent, inspectable integration of standard techniques with an
explicit feature contract and a safety-first reconciliation---reported against an honest
evidential standard rather than an inflated one. In answering its six research questions,
the work shows that such a pipeline can be built and made transparent, while remaining
candid that demonstrating real-world predictive validity lies outside its present scope.

\subsection{Future Work}\label{subsec:future-work}
The limitations map directly onto a program of future work, ordered by how much each would
strengthen the evidential standing of the system:

\begin{itemize}
\item[\textbf{1)}] Replace the clustering-derived proxy labels with recorded default or
  loss outcomes, so the classifier predicts a ground-truth target rather than a clustering
  artifact \cite{ref18,ref19}.
\item[\textbf{2)}] Adopt entity-level (grouped) splitting by identifier to remove the
  entity-leakage concern and yield a trustworthy generalization estimate.
\item[\textbf{3)}] Break the arithmetic label dependency by constructing labels from
  information not fully recoverable from the classifier inputs, or by explicitly modeling
  that dependency.
\item[\textbf{4)}] Rebalance the guardrail and the model---including resolving the 60/75
  medium-threshold inconsistency---so the learned component contributes to more decisions.
\item[\textbf{5)}] Strengthen the explanation layer by adopting SHAP or a partial-derivative
  sensitivity method and combining attributions, as the benchmark literature recommends
  \cite{ref10,ref21}.
\item[\textbf{6)}] Integrate historical or live market and DeFi data, modeling the
  liquidation and stablecoin-instability risks documented in the literature, to turn the
  deterministic projection into a genuine backtest \cite{ref5,ref6}.
\item[\textbf{7)}] Deploy a real audit ledger via a smart contract, with attention to
  storage, throughput, and oracle-trust costs, to make the auditability claim
  production-grade \cite{ref13,ref14,ref15,ref22}.
\item[\textbf{8)}] Move toward dynamic, causal allocation---reinforcement-learning
  rebalancing or an explainable optimizer---framed with the causal-inference caution raised
  for financial planning \cite{ref4,ref8,ref12}.
\end{itemize}

Taken together, these directions would progressively convert the present prototype---an
honest demonstration of integration---into a system whose predictive and financial claims
could be defended on external evidence. That progression, rather than any single algorithmic
advance, is the path these findings point toward.
